\documentclass[12pt]{article}
\usepackage{amsmath, amssymb}
\usepackage{geometry}
\usepackage{enumitem}
\geometry{margin=1in}

\begin{document}

\title{CSE400 – Fundamentals of Probability in Computing\\
Lecture 10: Randomized Min-Cut Algorithm}
\date{}
\maketitle

\section*{1. Min-Cut Problem}

\subsection*{1.1 Why Use Min-Cut?}

We use the min-cut algorithm in various applications to solve problems related to:

Network connectivity

Reliability

Optimization

Applications:

Network Design

Min-cut helps in improving the efficiency of communication and optimizing network flow.

The algorithm is used in network design to find the minimum capacity cut.

Communication Networks

To understand network vulnerability to failures.

Helps in building robust and fault-tolerant communication networks.

VLSI Design

Useful for partitioning circuits into smaller components.

Leads to reduced interconnectivity complexity.

(Reference: Section 1.5 – Application: A Randomized Min-Cut Algorithm in Probability and Computing (2nd Edition))

\section*{2. What is Min-Cut?}

\subsection*{2.1 Cut-Set Definition}

A cut-set in a graph is:

A set of edges whose removal breaks the graph into two or more connected components.

\subsection*{2.2 Minimum Cut Problem}

Given a graph

\[
G = (V, E)
\]

with \( n \) vertices,

The minimum cut (min-cut) problem is:

To find a minimum cardinality cut-set in \( G \).

\subsection*{2.3 Edge Contraction (Main Operation)}

The main operation in the algorithm is edge contraction.

Edge contraction:

Removes an edge from the graph

Merges the two vertices connected by that edge

If we contract edge \( (u, v) \):

Merge vertices \( u \) and \( v \) into one vertex

Eliminate all edges connecting \( u \) and \( v \)

Retain all other edges

After contraction:

Parallel edges may exist

No self-loops are allowed

\section*{3. Successful and Unsuccessful Min-Cut Run}

\subsection*{3.1 Successful Min-Cut Run}

A successful min-cut run refers to:

The success in the outcome of an algorithm designed to find the minimum cut in a graph.

\subsection*{3.2 Unsuccessful Min-Cut Run}

An unsuccessful min-cut run refers to:

An iteration of a min-cut algorithm where the algorithm fails to correctly identify the minimum cut of a given graph.

Min-cut algorithms (such as Karger’s algorithm) are random and sensitive to initial edge choices.

If critical edges are contracted early:

The algorithm may not find the correct minimum cut.

\section*{4. Max-Flow Min-Cut Theorem}

Theorem Statement

In a flow network, the maximum amount of flow passing from the source to the sink is equal to the total weight of the edges in a minimum cut.

Definitions

Capacity of a cut

Sum of capacities of edges oriented from a vertex in set \( X \) to a vertex in set \( Y \).

Minimum cut

The cut having the smallest possible capacity.

Minimum cut capacity

Capacity of the minimum cut.

Maximum flow

Largest possible flow from source \( S \) to sink \( T \).

\section*{5. Deterministic Min-Cut Algorithm}

Stoer–Wagner Min-Cut Algorithm

Let \( s \) and \( t \) be two vertices of graph \( G \).

Let \( G/\{s,t\} \) be the graph obtained by merging \( s \) and \( t \).

Then:

A minimum cut of \( G \) can be obtained by taking the smaller of:

A minimum \( s - t \) cut of \( G \)

A minimum cut of \( G/\{s,t\} \)

Reason:

If a minimum cut separates \( s \) and \( t \):

Then minimum \( s - t \) cut is the minimum cut.

If no minimum cut separates \( s \) and \( t \):

Then minimum cut of \( G/\{s,t\} \) is the answer.

\subsection*{5.1 Pseudocode}

Algorithm 1: MinimumCutPhase(G, a)

\( A \leftarrow \{a\}; \)

while \( A \neq V \) do

\hspace*{1cm} add to \( A \) the most tightly connected vertex;

return the cut weight as the "cut of the phase";

Algorithm 2: MinimumCut(G)

while \( |V| \geq 1 \) do

\hspace*{1cm} choose any \( a \) from \( V \);

\hspace*{1cm} MinimumCutPhase(G, a);

\hspace*{1cm} if cut-of-the-phase is lighter than current minimum cut then

\hspace*{2cm} store cut-of-the-phase as current minimum cut;

\hspace*{1cm} shrink \( G \) by merging the two vertices added last;

return the minimum cut;

Time Complexity

Stoer-Wagner algorithm:

\[
O(V \cdot E + V^2 \log V)
\]

Where:

\( V \) = number of vertices

\( E \) = number of edges

Deterministic min-cut:

Always guarantees exact minimum cut

May have higher time complexity for large graphs

\section*{6. Randomized Min-Cut Algorithm}

\subsection*{6.1 Why Randomized Algorithm?}

Randomized algorithms provide:

Probabilistic guarantee of success

More accurate estimate with fewer iterations

Advantages:

Efficiency

Parallelization

Approximation guarantees

Avoidance of worst-case instances

Heuristic nature

Robustness

\section*{7. Karger’s Randomized Algorithm}

(Uses repeated random edge contraction.)

The algorithm repeatedly:

Randomly selects an edge

Contracts it

Continues until only two vertices remain

The remaining edges between the two vertices form a cut.

\section*{8. Comparison: Deterministic vs Randomized Min-Cut}

Deterministic Min-Cut

Guarantees exact minimum cut

Higher time complexity

Time complexity:

\[
O(V \cdot E + V^2 \log V)
\]

Randomized Min-Cut (Karger’s Algorithm)

Provides approximate minimum cut with high probability

Time complexity:

\[
O(V^2)
\]

Choice of approach depends on the specific problem.

\section*{9. Theorem for Min-Cut Set}

The algorithm outputs a min-cut set with probability at least:

\[
\frac{2}{n(n-1)}
\]

Where:

\( n \) = number of vertices

\section*{10. Python Simulation}

Students were instructed to:

Open Campuswire post for Lecture 10

Download the provided .ipynb file

Run Python simulation

\section*{Final Summary (Exam Revision Points)}

Cut-set → edges whose removal disconnects graph

Min-cut → minimum cardinality cut-set

Edge contraction → merge two vertices

Max-flow = Min-cut capacity

Stoer-Wagner → deterministic, exact

Karger → randomized, probabilistic

Probability of success ≥ 
\[
\frac{2}{n(n-1)}
\]

Deterministic complexity:
\[
O(VE + V^2 \log V)
\]

Randomized complexity:
\[
O(V^2)
\]

\end{document}